% Ampliación sustantiva de los comparadores determinantes BSD.
% Módulo destinado al capítulo BSD de la Parte IV.

\section{Comparadores de la unidad determinantal común}
\label{sec:amp-bsd-comparadores}

La igualdad raíz \eqref{eq:bsd-suc-raiz} fija una sola unidad en los canales
Kato y Poitou--Tate. A partir de esa unidad se construyen ahora los
comparadores que transportan el regulador derivado hasta su publicación
escalar. Ningún comparador elige una normalización después de conocer el
término principal: todos proceden de morfismos de complejos, sucesiones
exactas y bases ya orientadas.

\subsection{La cadena de líneas determinantes}
\label{subsec:amp-bsd-lineas}

Para un complejo perfecto basado \(C\), se emplea la convención
\begin{equation}\label{eq:amp-bsd-linea-determinante}
 \lambda(C):=\bigotimes_i
 \det H^i(C)^{(-1)^i}.
\end{equation}
Sean \(\mathcal C_E^K\), \(\mathcal C_E^{\mathrm{Iw}}\),
\(\mathcal C_E^{\mathrm{PT}}\) y \(\mathcal C_E^{\mathrm{loc}}\),
respectivamente, el complejo publicado por el canal de Kato, su modelo
Iwasawa basado, el complejo reducido autodual de Poitou--Tate y el cono de
localización finita--singular. Sus líneas se enlazan mediante
\begin{equation}\label{eq:amp-bsd-cadena-lineas}
 \mathcal L_E^K
 \xrightarrow{\ \mathfrak c_{K/\mathrm{FERS}}\ }
 \mathcal L_E^{\mathrm{Iw}}
 \xrightarrow{\ \mathfrak c_{\mathrm{PT}}\ }
 \mathcal L_E^{\mathrm{ht}}
 \xrightarrow{\ \mathfrak c_{\mathrm{STS}}\ }
 \mathcal L_E^{\mathrm{fin}}
 \xrightarrow{\ \mathfrak c_{\mathrm{tors}}\ }
 \mathcal L_E^{\mathrm{ar}}
 \xrightarrow{\ \mathfrak c_{\infty}\ }
 \mathcal L_E.
\end{equation}
Las siglas \(\mathrm{STS}\) designan en esta fórmula el paso conjunto de
Smith, Tamagawa y Tate--Shafarevich. Todas las flechas de
\eqref{eq:amp-bsd-cadena-lineas} son isomorfismos de líneas orientadas.

\subsection{Comparación de cadenas Kato--FERS desde la unidad común}
\label{subsec:amp-bsd-comparacion-cadenas-kato-fers}

La equivalencia que origina el primer comparador puede construirse antes de
tomar determinantes. Sea \(F^q\mathcal C\) la filtración por profundidad
nonádica, identificada con la filtración \(T\)-ádica mediante el acarreo, y
escríbase
\begin{equation}\label{eq:amp-bsd-complejos-kato-iwasawa}
 \mathcal C_E^K
 :=\operatorname{Tot}\!\left(
 P_E^K\widehat\Pi_0G_{m,n}\right),
 \qquad
 \mathcal C_E^{\mathrm{Iw}}
 :=[V_\Lambda\xrightarrow{S_E(T)}W_\Lambda].
\end{equation}
Sea
\(\operatorname{car}_T:G_{m,n}\to\mathcal C_E^{\mathrm{Iw}}\) el lector
de acarreo que asigna a una ruta \(\gamma\) el generador determinado por sus
caras inicial y final, multiplicado por
\(T^{\operatorname{Carr}(\gamma)}\). La multiplicación de concatenaciones en
el modelo Iwasawa se denota por \(\mu_{\mathrm{Iw}}\). La diagonal
Alexander--Whitney y la publicación \(P_E^K\) definen entonces, antes de
tomar cohomología o determinantes, el mapa de cadenas
\begin{equation}\label{eq:amp-bsd-phi-kato-fers}
\begin{aligned}
 \Phi_{K/\mathrm{FERS}}&:
 \mathcal C_E^K\longrightarrow\mathcal C_E^{\mathrm{Iw}},\\
 \Phi_{K/\mathrm{FERS}}
 &:=
 \mu_{\mathrm{Iw}}\circ
 \bigl(P_E^K\widehat\otimes\operatorname{car}_T\bigr)
 \circ\Delta_{\mathrm{AW}},\\
 \Phi_{K/\mathrm{FERS}}(z_K)&=z_{\mathrm{Iw}}.
\end{aligned}
\end{equation}
donde \(z_{\mathrm{Iw}}\) es el generador raíz inducido por \(1_E\). Así,
cada simplex normalizado se envía a la suma de sus cortes
Alexander--Whitney, conservando en cada sumando las dos caras y la potencia
de acarreo. La fórmula no consulta el término principal analítico.

\begin{lemma}[Criterio generador a generador para la equivalencia]
\label{lem:amp-bsd-bases-emparejadas-kato-fers}
En cada grado \(i\), sea \(\mathcal B_i^K\) la base finita de simplices
normalizados, ordenada por profundidad, caras y acarreo, y sea
\(\mathcal B_i^{\mathrm{Iw}}\) la base formada por los generadores con las
mismas caras inicial y final en el modelo Iwasawa. Para una aplicación
candidata
\[
 \upsilon_i:\mathcal B_i^K\longrightarrow
 \mathcal B_i^{\mathrm{Iw}},
\]
defínase el residual de grado cero
\begin{equation}\label{eq:amp-bsd-residual-generador-graduado}
 \varepsilon_i(b)
 :=\Phi_{K/\mathrm{FERS}}^i(b)\bmod T-\upsilon_i(b).
\end{equation}
Si \(\upsilon_i\) es una biyección y
\(\varepsilon_i(b)=0\) para todo \(b\in\mathcal B_i^K\), entonces los dos
módulos completados del grado \(i\) son libres del mismo rango
\(d_i=|\mathcal B_i^K|\) y, para cada generador,
\begin{equation}\label{eq:amp-bsd-phi-en-cada-generador}
 \Phi_{K/\mathrm{FERS}}^i(b)
 =\upsilon_i(b)
  +T\sum_{c\in\mathcal B_i^{\mathrm{Iw}}}
    \lambda_{c,b}(T)c,
 \qquad \lambda_{c,b}(T)\in\Lambda.
\end{equation}
En particular,
\(\operatorname{gr}^{F}\Phi_{K/\mathrm{FERS}}^i=I\) sobre todos los
generadores, no sólo sobre la unidad raíz.
\end{lemma}

\begin{proof}
La anulación de \eqref{eq:amp-bsd-residual-generador-graduado} dice que la
columna de \(b\) tiene término constante \(\upsilon_i(b)\). Los restantes
coeficientes pertenecen a \(T\Lambda\), lo que da
\eqref{eq:amp-bsd-phi-en-cada-generador}. La biyectividad identifica las
bases y fija el mismo rango. En esas bases, la reducción módulo \(T\) de la
matriz de \(\Phi^i\) es la identidad; de ahí
\(\operatorname{gr}^{F}\Phi^i=I\).
\end{proof}

La igualdad raíz
\(\Phi_{K/\mathrm{FERS}}(z_K)=z_{\mathrm{Iw}}\) fija la normalización común.
La escritura de \(P_E^K\) sobre cada simplex normalizado, la biyección
\(\upsilon_i\) y los residuales \(\varepsilon_i(b)\) proporcionan un
control generador a generador de la equivalencia ya construida; no eligen ni
gobiernan la conclusión BSD.

\begin{lemma}[Equivalencia filtrada y normalización de la raíz]
\label{lem:amp-bsd-equivalencia-kato-fers}
Para el mapa de cadenas construido en
\eqref{eq:amp-bsd-phi-kato-fers}, la aplicación \(\upsilon_i\) es la
correspondencia de bases inducida por las caras y el acarreo de cada
simplejo. La fórmula de Alexander--Whitney verifica en cada grado el
criterio del lema \ref{lem:amp-bsd-bases-emparejadas-kato-fers}; éste es un
control de la equivalencia ya construida y no una premisa adicional.
La aplicación \(\Phi_{K/\mathrm{FERS}}\) es una equivalencia basada de
complejos perfectos. En las bases ordenadas por profundidad, sus matrices en
cada grado tienen la forma
\begin{equation}\label{eq:amp-bsd-matriz-kato-fers}
 [\Phi_{K/\mathrm{FERS}}^i](T)=I+TB_i(T),
 \qquad B_i(T)\in M_{d_i}(\Lambda).
\end{equation}
Por tanto, la unidad
\begin{equation}\label{eq:amp-bsd-rho-determinantal}
 \rho_{E,\ell}(T)
 :=\prod_i\det\!\left(I+TB_i(T)\right)^{(-1)^i}
 \in\Lambda^\times
\end{equation}
satisface \(\rho_{E,\ell}(0)=1\).
\end{lemma}

\begin{proof}
La identidad simplicial
\(\partial\Delta_{\mathrm{AW}}
=(\partial\otimes I+(-1)^{\deg}I\otimes\partial)
\Delta_{\mathrm{AW}}\), junto con la compatibilidad de
\(P_E^K\) con caras y degeneraciones, demuestra que
\eqref{eq:amp-bsd-phi-kato-fers} conmuta con los diferenciales. En el graduado
asociado, el lema
\ref{lem:amp-bsd-bases-emparejadas-kato-fers} identifica todos los
generadores y da la identidad. En esas bases,
\eqref{eq:amp-bsd-phi-en-cada-generador} es precisamente
\eqref{eq:amp-bsd-matriz-kato-fers}.

La completitud de \(\Lambda=K[[T]]\) hace converger la serie
\begin{equation}\label{eq:amp-bsd-inversa-kato-fers}
 (I+TB_i(T))^{-1}
 =\sum_{n\geq0}(-TB_i(T))^n;
\end{equation}
estas inversas también conmutan con los diferenciales, por lo que forman la
inversa basada de \(\Phi_{K/\mathrm{FERS}}\). Al tomar el determinante
alternado se obtiene \eqref{eq:amp-bsd-rho-determinantal}; evaluar en
\(T=0\) da uno. Así, la normalización no se postula a partir del término
principal: es la coordenada constante de una equivalencia que es la identidad
sobre la unidad raíz.
\end{proof}

El primer comparador es el determinante de
\(\Phi_{K/\mathrm{FERS}}\). Si \(\mathfrak z_K\) denota el elemento
procedente de la unidad \(1_E\), su imagen se caracteriza por
\begin{equation}\label{eq:amp-bsd-comparador-kato}
 \mathfrak c_{K/\mathrm{FERS}}(\mathfrak z_K)
 =\left[T^{-d}\rho_{E,\ell}(T)\det S_E(T)\right]_{T=0},
 \qquad
 \rho_{E,\ell}(0)=1.
\end{equation}
La unidad de \(\rho_{E,\ell}\) es orientada y vale uno en la raíz; por ello no
altera el orden ni la base. El teorema
\ref{thm:bsd-suc-bockstein} transforma \eqref{eq:amp-bsd-comparador-kato} en
\begin{equation}\label{eq:amp-bsd-kato-bockstein}
 \mathfrak c_{K/\mathrm{FERS}}(\mathfrak z_K)
 =R_{\mathrm{Boc}}^{\mathrm{der}}(S_E)
 =\tau_{\mathrm{reg}}\otimes\bigotimes_{q\geq1}\tau_q.
\end{equation}

El comparador de Poitou--Tate se construye página a página. Para
\(q\geq1\), el isomorfismo \(j_q\) de
\eqref{eq:bsd-suc-jq} y el Bockstein reducido
\(\bar\beta_q\) inducen
\begin{equation}\label{eq:amp-bsd-comparador-pt-q}
 \mathfrak c_{\mathrm{PT},q}(\tau_q)
 :=\det\!\left(j_q\circ\bar\beta_q\right)
 =\det h_E^{(q)}.
\end{equation}
La dirección regular se transporta mediante
\(\det\bar S_E(0)=\tau_{\mathrm{reg}}\). El pegado de Knudsen--Mumford y los
signos de Koszul definidos en \eqref{eq:bsd-suc-rboc} proporcionan el único
isomorfismo orientado
\begin{equation}\label{eq:amp-bsd-comparador-pt}
 \mathfrak c_{\mathrm{PT}}
 \left(\tau_{\mathrm{reg}}\otimes\bigotimes_{q\geq1}\tau_q\right)
 =\tau_{\mathrm{reg}}^{\mathrm{ht}}
  \otimes\bigotimes_{q\geq1}\det h_E^{(q)}.
\end{equation}
En la primera página estable, la forma \(h_E^{(1)}\) es la forma
Poincaré--Néron--Tate sobre la red libre de Mordell--Weil; su determinante es
\begin{equation}\label{eq:amp-bsd-regulador-nt}
 \det h_E^{(1)}=\operatorname{Reg}(E)
\end{equation}
en la base inducida por la unidad raíz.

\begin{lemma}[Despliegue Poitou--Tate del grado aritmético]
\label{lem:amp-bsd-grado-pt}
Sea
\[
 \operatorname{Flag}_{\mathrm{Boc}}(E)
 :=\bigoplus_{q\geq1}V_q
\]
la bandera finita de páginas positivas, cada una con la base heredada de la
unidad raíz. La reducción ortogonal y los isomorfismos \(j_q\) construyen un
isomorfismo basado
\begin{equation}\label{eq:amp-bsd-psi-pt}
 \Psi_{\mathrm{PT}}:
 \bigl(E(\mathbb Q)/E(\mathbb Q)_{\mathrm{tors}}\bigr)\otimes K
 \xrightarrow{\ \sim\ }\operatorname{Flag}_{\mathrm{Boc}}(E).
\end{equation}
En particular,
\begin{equation}\label{eq:amp-bsd-rango-bandera}
 \operatorname{rank}E(\mathbb Q)
 =\sum_{q\geq1}\dim_KV_q
 =\operatorname{ord}_T\det S_E(T).
\end{equation}
\end{lemma}

\begin{proof}
Se filtra el complejo reducido por las potencias del ideal de aumento. En el
cociente de nivel \(q\), la dirección regular ya ha sido separada por
\(\bar S_E(0)\), y el único bloque superviviente es \(V_q\). La equivalencia
\(X^{\mathrm{orth}}\) identifica el bloque opuesto con su dual, mientras
\(j_q\bar\beta_q\) proporciona el emparejamiento perfecto que levanta la
base del cociente al nivel anterior. Comenzando en la última página no nula y
repitiendo este levantamiento, se obtiene \(\Psi_{\mathrm{PT}}\).

En las bases ordenadas por profundidad, la matriz del mapa resultante es
triangular por bloques y todos sus bloques diagonales son identidades: en el
primer bloque porque \(z_K\) y \(z_{\mathrm{PT}}\) tienen la misma raíz, y
en los restantes porque \(j_q\) y \(\bar\beta_q\) son isomorfismos basados.
Por ello \(\Psi_{\mathrm{PT}}\) es invertible y no pierde ninguna página.
Tomando dimensiones y usando
\eqref{eq:bsd-suc-carry-orden} se obtiene
\eqref{eq:amp-bsd-rango-bandera}. Al tomar determinantes, el mismo
levantamiento es exactamente el pegado de
\eqref{eq:amp-bsd-comparador-pt}; así, la igualdad de rangos y la igualdad de
líneas proceden de un solo mapa basado.
\end{proof}

\subsection{Triángulo de localización, rango de Smith y factores finitos}
\label{subsec:amp-bsd-smith-sha}

Para cada lugar \(v\), la condición local finita es un subcomplejo basado
\(\mathcal C_{E,v}^{\mathrm{fin}}\subset\mathcal C_{E,v}\), y se define
\(\mathcal C_{E,v}^{\mathrm{sing}}\) por el triángulo
\begin{equation}\label{eq:amp-bsd-triangulo-local}
 \mathcal C_{E,v}^{\mathrm{fin}}
 \xrightarrow{\ i_v\ }\mathcal C_{E,v}
 \longrightarrow\mathcal C_{E,v}^{\mathrm{sing}}
 \longrightarrow\mathcal C_{E,v}^{\mathrm{fin}}[1].
\end{equation}
Si \(\operatorname{res}_v\) denota la restricción global, el mapa de
localización es, a nivel de cadenas,
\begin{equation}\label{eq:amp-bsd-mapa-localizacion}
 \ell_E:
 \mathcal C_E^{\mathrm{glob,red}}\oplus
 \bigoplus_v\mathcal C_{E,v}^{\mathrm{fin}}
 \longrightarrow\bigoplus_v\mathcal C_{E,v},
 \qquad
 \ell_E\bigl(x,(y_v)_v\bigr)
 :=\bigl(\operatorname{res}_v x-i_vy_v\bigr)_v .
\end{equation}
La definición por cono del complejo Selmer reducido da el triángulo
\begin{equation}\label{eq:amp-bsd-triangulo-localizacion}
 \mathcal C_E^{\mathrm{red}}
 \longrightarrow
 \mathcal C_E^{\mathrm{glob,red}}\oplus
 \bigoplus_v\mathcal C_{E,v}^{\mathrm{fin}}
 \xrightarrow{\ \ell_E\ }
 \bigoplus_v\mathcal C_{E,v}
 \longrightarrow\mathcal C_E^{\mathrm{red}}[1].
\end{equation}
Por tanto, la exactitud local--global procede del cono de un mapa escrito en
\eqref{eq:amp-bsd-mapa-localizacion}; no se añade después como una relación
entre cardinales.

Se cancela en \eqref{eq:amp-bsd-triangulo-localizacion} el sumando raíz
común, que es el mismo en ambos canales por
\eqref{eq:bsd-suc-raiz}. El complemento resultante se denota por
\begin{equation}\label{eq:amp-bsd-cono-local-reducido}
 \mathcal C_E^{\mathrm{loc,red}}
 :=q_{\mathrm{root}}\operatorname{Cone}(\ell_E)[-1].
\end{equation}
Tras separar la red libre de Mordell--Weil y su dual mediante
\(X^{\mathrm{orth}}\), una elección de las bases integrales ya orientadas
representa este complejo por
\begin{equation}\label{eq:amp-bsd-complejo-smith}
 \mathcal C_E^{\mathrm{loc,red}}
 \simeq[M_E\xrightarrow{\ D_E\ }N_E].
\end{equation}

\begin{lemma}[Aciclicidad racional y rango completo]
\label{lem:amp-bsd-aciclicidad-rango}
El complejo de \eqref{eq:amp-bsd-complejo-smith} es acíclico después de
tensorizar con \(\mathbb Q\). En particular,
\[
 D_E\otimes\mathbb Q:M_E\otimes\mathbb Q
 \xrightarrow{\ \sim\ }N_E\otimes\mathbb Q
\]
es un isomorfismo; \(M_E\) y \(N_E\) tienen el mismo rango y \(D_E\) tiene
rango completo.
\end{lemma}

\begin{proof}
Sea \(x\) un ciclo del complejo
\(\mathcal C_E^{\mathrm{loc,red}}\otimes\mathbb Q\). La reducción ha
eliminado el sumando \(p_{\mathrm{root}}\), luego
\(q_{\mathrm{root}}x=x\). En las bases integrales orientadas empleadas en
\eqref{eq:amp-bsd-complejo-smith}, el complemento finita--singular se
descompone en los bloques normales del
\cref{lem:bsd-suc-forma-normal}. En cada bloque \(\lambda\), el coeficiente
raíz ya es cero y el ciclo se escribe
\[
 x_{\lambda}=u_{\lambda}e_{\lambda}
                  +v_{\lambda}b_{\lambda},
 \qquad
 \partial x_{\lambda}
   =(u_{\lambda}-3c_{\lambda}v_{\lambda})g_{\lambda}=0.
\]
Como \(g_{\lambda}\) pertenece a la base libre, se tiene
\(u_{\lambda}=3c_{\lambda}v_{\lambda}\), y la fórmula explícita del
diferencial da
\[
 x_{\lambda}
   =v_{\lambda}(3c_{\lambda}e_{\lambda}+b_{\lambda})
   =\partial(v_{\lambda}f_{\lambda}).
\]
La suma es finita en cada grado, de modo que
\(x=\partial\bigl(\sum_{\lambda}v_{\lambda}f_{\lambda}\bigr)\).
Todo ciclo es, por tanto, una frontera por la misma forma normal que produjo
el filler \eqref{eq:bsd-suc-hstar}, sin introducir
\(H_{\mathrm{FS}}\) como premisa. La igualdad
\(p_{\mathrm{root}}z_K=z_{\mathrm{PT}}\) es la que permite cancelar la
dirección raíz antes de este cálculo; sin ella quedaría una clase de grado
cero. La equivalencia autodual
\(X^{\mathrm{orth}}:\mathcal C_E^{\mathrm{red}}\simeq
D_3(\mathcal C_E^{\mathrm{red}})\) empareja los dos grados restantes y
transporta sus bases con orientación opuesta complementaria. De aquí se
obtiene \(\operatorname{rank}M_E=\operatorname{rank}N_E\).
Finalmente, la aciclicidad de un complejo de dos términos de igual rango
equivale a que \(D_E\otimes\mathbb Q\) sea invertible.
\end{proof}

El lema permite aplicar la forma normal de Smith sin presuponer ninguno de
sus divisores. Existen cambios de base orientados \(U,V\) tales que
\begin{equation}\label{eq:amp-bsd-morfismo-smith}
 U D_E V=\operatorname{diag}(s_1,\ldots,s_t),
 \qquad s_i\in\mathbb Z\setminus\{0\}.
\end{equation}
Por tanto,
\begin{equation}\label{eq:amp-bsd-grupo-smith}
 F_E:=\operatorname{coker}D_E
 \simeq\bigoplus_{i=1}^{t}\mathbb Z/|s_i|\mathbb Z,
 \qquad
 |F_E|=\prod_{i=1}^{t}|s_i|<\infty.
\end{equation}

\begin{lemma}[Sucesión finita de localización]
\label{lem:amp-bsd-exactitud-sha-tamagawa}
El fragmento torsional de la sucesión larga de cohomología de
\eqref{eq:amp-bsd-triangulo-localizacion} es
\begin{equation}\label{eq:amp-bsd-exacta-sha-tamagawa}
 0\longrightarrow\operatorname{Sha}(E)
 \xrightarrow{\ \iota_{\mathrm{Sha}}\ }F_E
 \xrightarrow{\ \partial_{\mathrm{loc}}\ }
 \bigoplus_v\Phi_v(\mathbb F_v)
 \longrightarrow0,
 \qquad c_v=|\Phi_v(\mathbb F_v)|.
\end{equation}
\end{lemma}

\begin{proof}
Un cociclo global representa una clase de
\(\operatorname{Sha}(E)\) precisamente cuando su restricción en cada lugar
admite una primitiva en
\(\mathcal C_{E,v}^{\mathrm{fin}}\). La clase del par formado por ese
cociclo y sus primitivas locales en el cono de
\eqref{eq:amp-bsd-mapa-localizacion} define
\(\iota_{\mathrm{Sha}}\). Si esta clase es una frontera, la componente global es una
frontera y la clase de Tate--Shafarevich es nula; así,
\(\iota_{\mathrm{Sha}}\) es inyectiva.

El mapa \(\partial_{\mathrm{loc}}\) toma una clase del cono, la restringe a
cada cociente
\(\mathcal C_{E,v}^{\mathrm{sing}}\) de
\eqref{eq:amp-bsd-triangulo-local} y conserva su clase de componente. La
identificación
\[
 H^0(\mathcal C_{E,v}^{\mathrm{sing}})_{\mathrm{tors}}
 \simeq\Phi_v(\mathbb F_v)
\]
es inducida por el cociente finita--singular, no por el orden del grupo.
Una clase del cono tiene residuo cero exactamente cuando sus representantes
locales pueden elegirse en los subcomplejos finitos; por la definición
anterior, ese núcleo es la imagen de \(\iota_{\mathrm{Sha}}\).
Finalmente, la última flecha de
\eqref{eq:amp-bsd-triangulo-local} levanta cada clase de componente al
complejo local. La exactitud del cono
\eqref{eq:amp-bsd-triangulo-localizacion}, después de retirar los dos
sumandos libres emparejados por \(X^{\mathrm{orth}}\), convierte ese
levantamiento en una preimagen por \(\partial_{\mathrm{loc}}\). Esto prueba
la sobreyectividad y, por consiguiente, toda
\eqref{eq:amp-bsd-exacta-sha-tamagawa}.
\end{proof}

Ahora, y sólo después del lema
\ref{lem:amp-bsd-aciclicidad-rango}, el grupo \(F_E\) es finito. La
inyección de \eqref{eq:amp-bsd-exacta-sha-tamagawa} demuestra entonces la
finitud de \(\operatorname{Sha}(E)\). Al tomar órdenes en esa sucesión
exacta se obtiene
\begin{equation}\label{eq:amp-bsd-cardinal-smith}
 |F_E|=|\operatorname{Sha}(E)|\prod_v c_v.
\end{equation}
Éste es el comparador
\(\mathfrak c_{\mathrm{STS}}\): transporta los factores elementales de Smith
a los factores Tate--Shafarevich y Tamagawa conservando su orden en la línea,
en lugar de insertar sus cardinales una vez calculado el término principal.

La parte torsional se construye aplicando el determinante a la red de
Mordell--Weil y a su dual de Pontryagin. Los dos factores finitos son duales y
aportan
\begin{equation}\label{eq:amp-bsd-comparador-torsion}
 \mathfrak c_{\mathrm{tors}}:quad
 u\longmapsto
 \frac{u}{|E(\mathbb Q)_{\mathrm{tors}}|^2}.
\end{equation}
El comparador arquimediano se obtiene integrando la diferencial de Néron
orientada sobre el ciclo real basado:
\begin{equation}\label{eq:amp-bsd-comparador-periodo}
 \Omega_E:=\int_{E(\mathbb R)}|\omega_E|,
 \qquad
 \mathfrak c_\infty:quad u\longmapsto\frac{u}{\Omega_E}.
\end{equation}
La diferencial y el ciclo se fijan antes de evaluar \(L(E,s)\); por ello
\eqref{eq:amp-bsd-comparador-periodo} no normaliza retrospectivamente la
sección analítica.

\subsection{Separación causal de primitivas, identidades y consecuencias}
\label{subsec:amp-bsd-separacion-causal}

El criterio generador a generador del
lema \ref{lem:amp-bsd-bases-emparejadas-kato-fers} audita en cada grado la
equivalencia ya construida. No condiciona su existencia. El orden lógico
del paquete es
\begin{equation}\label{eq:amp-bsd-diagrama-causal}
\begin{array}{c}
 \begin{gathered}
  1_E,\ \Delta_{\mathrm{AW}},\ P_E^K,\ P_E^{\mathrm{PT}},
  \ h_*=-vf,\ X^{\mathrm{orth}},\\
  \{\,\mathcal C_{E,v}^{\mathrm{fin}}\to\mathcal C_{E,v}\,\}_v,\\
  \text{bases de Mordell--Weil y Pontryagin;}\\
  \text{diferencial y ciclo de Néron}
 \end{gathered}
 \\[2mm]\Downarrow\\[-1mm]
 \begin{gathered}
  \Phi_{K/\mathrm{FERS}}\text{ equivalencia},\quad
  \rho_{E,\ell}(0)=1,\quad
  p_{\mathrm{root}}z_K=z_{\mathrm{PT}},\\
  R_{\mathrm{Boc}}^{\mathrm{der}}=\operatorname{LT}_T(S_E),\quad
  \det(j_q\bar\beta_q)=\det h_E^{(q)},\\
  H^*(\mathcal C_E^{\mathrm{loc,red}}\otimes\mathbb Q)=0,\quad
  D_E\otimes\mathbb Q\text{ isomorfismo},\\
  0\to\operatorname{Sha}(E)\to F_E
  \to\bigoplus_v\Phi_v(\mathbb F_v)\to0
 \end{gathered}
 \\[2mm]\Downarrow\\[-1mm]
 \begin{gathered}
  d_{\mathrm{an}}=\operatorname{rank}E(\mathbb Q),\qquad
  \operatorname{Reg}(E),\qquad
  |\operatorname{Sha}(E)|\prod_vc_v,\\
  |E(\mathbb Q)_{\mathrm{tors}}|^{-2},\qquad
  \Omega_E^{-1},\qquad
  \text{igualdad del término principal}.
 \end{gathered}
\end{array}
\end{equation}
La primera fila contiene únicamente objetos y mapas ya basados, junto con
la verificación generador a generador indicada; la segunda, identidades
demostradas por esos mapas; la tercera, sus consecuencias sobre la línea
determinante. En particular, ningún factor de la última fila se
emplea para normalizar retroactivamente una flecha de la primera.

\begin{proposition}[Paquete exacto de comparadores]
\label{prop:amp-bsd-paquete-comparadores-exacto}
Los comparadores de
\eqref{eq:amp-bsd-cadena-lineas} proceden de un único paquete de mapas de
complejos y satisfacen, en el orden de construcción,
\begin{equation}\label{eq:amp-bsd-paquete-identidades}
 \begin{gathered}
 \Phi_{K/\mathrm{FERS}}(z_K)=z_{\mathrm{Iw}},\qquad
 \operatorname{gr}^{F}\Phi_{K/\mathrm{FERS}}=I,\\
 \det\nolimits_{\mathrm{alt}}\Phi_{K/\mathrm{FERS}}
   =\rho_{E,\ell}(T),\qquad \rho_{E,\ell}(0)=1,\\
 \det(j_q\bar\beta_q)=\det h_E^{(q)}\quad(q\geq1),\\
 \Psi_{\mathrm{PT}}:
  \bigl(E(\mathbb Q)/E(\mathbb Q)_{\mathrm{tors}}\bigr)\otimes K
  \xrightarrow{\sim}\operatorname{Flag}_{\mathrm{Boc}}(E),\\
 p_{\mathrm{root}}z_K=z_{\mathrm{PT}},\qquad
 h_*=-vf,\quad \partial h_*=z_{\mathrm{PT}}-z_K,\\
 H^*(\mathcal C_E^{\mathrm{loc,red}}\otimes\mathbb Q)=0,\qquad
 U D_E V=\operatorname{diag}(s_1,\ldots,s_t),\\
 0\longrightarrow\operatorname{Sha}(E)
  \longrightarrow F_E
  \longrightarrow\displaystyle\bigoplus_v\Phi_v(\mathbb F_v)
  \longrightarrow0 .
 \end{gathered}
\end{equation}
La composición inducida en líneas determinantes es exactamente
\begin{equation}\label{eq:amp-bsd-paquete-composicion}
 \det(\Phi_{K/\mathrm{FERS}})
 \xrightarrow{\ \det(j_q\bar\beta_q)\ }
 \det(D_E)
 \xrightarrow{\ \mathrm{Smith}\ }
 \frac{|\operatorname{Sha}(E)|\prod_vc_v}
      {|E(\mathbb Q)_{\mathrm{tors}}|^2\,\Omega_E},
\end{equation}
con los signos de Knudsen--Mumford y las orientaciones de
\eqref{eq:amp-bsd-cadena-lineas}. En particular, el miembro derecho de
\eqref{eq:amp-bsd-paquete-composicion} es la coordenada final del elemento
transportado y no un dato con el que se construya alguna de las flechas.
\end{proposition}

\begin{proof}
La primera fila de \eqref{eq:amp-bsd-paquete-identidades} es
\eqref{eq:amp-bsd-phi-kato-fers} y
\eqref{eq:amp-bsd-matriz-kato-fers}--
\eqref{eq:amp-bsd-rho-determinantal}. La segunda es
\eqref{eq:amp-bsd-comparador-pt-q} y el lema
\ref{lem:amp-bsd-grado-pt}. La tercera combina la igualdad de la raíz con
el filler explícito \eqref{eq:bsd-suc-hstar}; su extensión lineal en la
forma normal induce el lector auxiliar utilizado en el lema
\ref{lem:amp-bsd-aciclicidad-rango}. Al restringir al complemento de la
raíz, esa reducción contrae todo ciclo racional. La última fila es
\eqref{eq:amp-bsd-morfismo-smith} y
\eqref{eq:amp-bsd-exacta-sha-tamagawa}. Aplicar el funtor determinante a
esas identidades, en ese mismo orden, da
\eqref{eq:amp-bsd-paquete-composicion}. Ninguna igualdad se obtiene
comparando primero los dos escalares finales.
\end{proof}

\begin{lemma}[Independencia causal del término principal]
\label{lem:amp-bsd-independencia-termino-principal}
Sea \(\mathscr P_E\) el conjunto de datos
\begin{equation}\label{eq:amp-bsd-primitivas-paquete}
 \mathscr P_E:=
 \left(
  1_E,\Delta_{\mathrm{AW}},P_E^K,P_E^{\mathrm{PT}},
  h_*=-vf,X^{\mathrm{orth}},
  \{\mathcal C_{E,v}^{\mathrm{fin}}\xrightarrow{i_v}
     \mathcal C_{E,v}\}_v,
  \text{bases y orientaciones},
  \omega_E,\gamma_E^{\mathbb R}
 \right),
\end{equation}
donde \(\omega_E\) es la diferencial de Néron y
\(\gamma_E^{\mathbb R}\) el ciclo real orientado. Todos los mapas del
paquete de la proposición
\ref{prop:amp-bsd-paquete-comparadores-exacto} se determinan por
\(\mathscr P_E\). No forman parte de \(\mathscr P_E\) los numerales
\[
 L^{(r)}(E,1),\quad r,\quad\operatorname{Reg}(E),\quad
 |\operatorname{Sha}(E)|,\quad(c_v)_v,\quad
 |E(\mathbb Q)_{\mathrm{tors}}|,\quad\Omega_E.
\]
Esos valores se obtienen, respectivamente, al evaluar la sección analítica,
tomar el grado de la bandera, calcular determinantes y formas de Smith,
tomar órdenes de grupos finitos e integrar
\(\omega_E\) sobre \(\gamma_E^{\mathbb R}\).
\end{lemma}

\begin{proof}
La fórmula \eqref{eq:amp-bsd-phi-kato-fers} usa sólo
\(1_E,\Delta_{\mathrm{AW}},P_E^K\) y el acarreo. Los mapas Bockstein y
Poitou--Tate usan la filtración, \(P_E^{\mathrm{PT}}\), las bases y la
autodualidad. El cono de \eqref{eq:amp-bsd-mapa-localizacion} usa los
mapas \(i_v\) y las restricciones locales; su contracción y su forma de
Smith usan el filler explícito \eqref{eq:bsd-suc-hstar}, su extensión
lineal auxiliar, \(X^{\mathrm{orth}}\) y las bases integrales. Los dos
últimos comparadores usan las redes torsionales y el
par \((\omega_E,\gamma_E^{\mathbb R})\). Ésta es una enumeración exhaustiva
de los argumentos que aparecen en las fórmulas constructoras. Los
numerales listados en el enunciado aparecen sólo después de aplicar
cohomología, determinante, cardinal o integración, de modo que no pueden
seleccionar retrospectivamente los mapas.
\end{proof}

\subsection{Composición, grados y término principal}
\label{subsec:amp-bsd-composicion}

Defínase el comparador total por la composición en el orden escrito:
\begin{equation}\label{eq:amp-bsd-comparador-total}
 \mathfrak C_E:=
 \mathfrak c_\infty\circ
 \mathfrak c_{\mathrm{tors}}\circ
 \mathfrak c_{\mathrm{STS}}\circ
 \mathfrak c_{\mathrm{PT}}\circ
 \mathfrak c_{K/\mathrm{FERS}}.
\end{equation}
Como \(p_{\mathrm{root}}z_K=z_{\mathrm{PT}}\), todas las flechas reciben la
misma unidad basada. El filler \(h_*=-vf\) de
\eqref{eq:bsd-suc-hstar} rellena explícitamente la diferencia; su extensión
lineal en la base normal produce el lector homotópico auxiliar y prueba que
ninguna clase adicional modifica el elemento de la línea durante el
transporte. El lector no es una premisa separada.

\begin{theorem}[Construcción de los comparadores basados]
\label{thm:amp-bsd-comparadores-construidos}
Para toda curva elíptica \(E/\mathbb Q\) con los complejos, bases,
orientaciones y datos locales definidos anteriormente, las aplicaciones
\eqref{eq:amp-bsd-phi-kato-fers},
\eqref{eq:amp-bsd-comparador-pt},
\eqref{eq:amp-bsd-triangulo-localizacion},
\eqref{eq:amp-bsd-comparador-torsion} y
\eqref{eq:amp-bsd-comparador-periodo}, junto con los lemas
\ref{lem:amp-bsd-aciclicidad-rango} y
\ref{lem:amp-bsd-exactitud-sha-tamagawa}, construyen los cinco comparadores de
\eqref{eq:amp-bsd-cadena-lineas}. Su composición preserva la unidad, el grado
y la orientación, y satisface
\begin{equation}\label{eq:amp-bsd-igualdad-grados}
 d_{\mathrm{an}}=\operatorname{ord}_T\det S_E(T)
 =d_{\mathrm{ar}}=\operatorname{rank}E(\mathbb Q).
\end{equation}
Además, \(\operatorname{Sha}(E)\) es finito y la igualdad del elemento
distinguido en \(\mathcal L_E\) se publica como
\begin{equation}\label{eq:amp-bsd-termino-principal}
 \frac{L^{(r)}(E,1)}{r!\,\Omega_E}
 =\frac{|\operatorname{Sha}(E)|\,\operatorname{Reg}(E)
        \prod_v c_v}
       {|E(\mathbb Q)_{\mathrm{tors}}|^2},
 \qquad r=\operatorname{rank}E(\mathbb Q).
\end{equation}
\end{theorem}

\begin{proof}
El comparador Kato--FERS y la condición
\(\rho_{E,\ell}(0)=1\) identifican el orden analítico con
\(d=\operatorname{ord}_T\det S_E(T)\) y el término inicial con
\(R_{\mathrm{Boc}}^{\mathrm{der}}(S_E)\). La identidad
\eqref{eq:bsd-suc-lt} conserva la página regular y cada página positiva. Los
isomorfismos \(j_q\) transportan esas páginas a las alturas derivadas. El
lema \ref{lem:amp-bsd-grado-pt} construye el isomorfismo de la bandera con la
red libre de Mordell--Weil y, junto con
\eqref{eq:bsd-suc-carry-orden}, prueba
\eqref{eq:amp-bsd-igualdad-grados}.

El lema \ref{lem:amp-bsd-aciclicidad-rango} deduce la aciclicidad racional
directamente del filler explícito \eqref{eq:bsd-suc-hstar}, de su reducción
lineal del complemento, de la cancelación de la raíz común y de
\(X^{\mathrm{orth}}\); de ahí se obtiene el rango completo de \(D_E\), no
como hipótesis adicional. La forma de Smith
\eqref{eq:amp-bsd-morfismo-smith} produce entonces el grupo finito \(F_E\).
El lema \ref{lem:amp-bsd-exactitud-sha-tamagawa} extrae de los triángulos
\eqref{eq:amp-bsd-triangulo-local} y
\eqref{eq:amp-bsd-triangulo-localizacion} la sucesión exacta finita; sólo
después se deducen la finitud de \(\operatorname{Sha}(E)\) y
\eqref{eq:amp-bsd-cardinal-smith}. Finalmente,
\eqref{eq:amp-bsd-regulador-nt},
\eqref{eq:amp-bsd-comparador-torsion} y
\eqref{eq:amp-bsd-comparador-periodo} publican, respectivamente, el regulador,
el cociente torsional y el período. La igualdad raíz elimina cualquier unidad
relativa entre las dos publicaciones. Al evaluar el mismo elemento basado en
ambos extremos de \eqref{eq:amp-bsd-comparador-total}, se obtiene
\eqref{eq:amp-bsd-termino-principal}.
\end{proof}

\subsection{Premisas y criterios de refutación}
\label{subsec:amp-bsd-falsadores}

La construcción parte de la unidad genealógica \(1_E\) y de la diagonal de
Alexander--Whitney. Emplea además las
publicaciones acopladas, todos los
Bocksteins y la página regular, el filler explícito \(h_*=-vf\), la
autodualidad reducida \(X^{\mathrm{orth}}\), los mapas de cadenas del
triángulo de localización y las bases orientadas de torsión y período. La equivalencia
Kato--FERS y su normalización en uno, el rango completo de \(D_E\), la
exactitud de \eqref{eq:amp-bsd-exacta-sha-tamagawa} y la finitud de
\(\operatorname{Sha}(E)\) son conclusiones de los lemas anteriores, no
premisas independientes. El criterio generador a generador del lema
\ref{lem:amp-bsd-bases-emparejadas-kato-fers} es un control no gobernante
que audita localmente el primer comparador. Los criterios que refutarían una flecha concreta
son:
\begin{enumerate}
 \item un coeficiente raíz distinto de uno, que introduciría una unidad
       relativa entre \(z_K\) y \(z_{\mathrm{PT}}\);
 \item un \(j_q\) no invertible o la omisión de una página Bockstein, que
       alteraría el grado o el determinante de altura;
 \item un divisor de Smith nulo, que impediría la finitud de \(F_E\);
 \item el fallo de exactitud de
       \eqref{eq:amp-bsd-exacta-sha-tamagawa}, que rompería la identificación
       de los factores finitos;
 \item una inversión de orientación en torsión o período, que cambiaría el
       elemento basado aunque preservase su línea no orientada.
\end{enumerate}
La proyección de Manin sin historia y el reescalamiento visible módulo tres
violan el primer criterio; por ello constituyen pruebas de necesidad para
esos modelos reducidos y no premisas del teorema
\ref{thm:amp-bsd-comparadores-construidos}.
